\section{Exact prime-cycle census}
\label{sec:census}

The quotient states make the recurrent computation much longer than a
trial-state count.  The exact length is the bridge from arithmetic execution
to operator compactness.

\begin{theorem}[Contracted expanded-$Q$ cycle length]
\label{thm:length}
Let $p$ be prime and $m=\lfloor\sqrt p\rfloor$.  The simple cycle
$\Gamma_p$ in \cref{prop:recurrent-census} has graph length
\begin{equation}
\boxed{
\ell(p)=2+\sum_{d=2}^{m}\left\lceil\frac pd\right\rceil
=2+\sum_{d=2}^{m}\left(1+\left\lfloor\frac pd\right\rfloor\right).}
\label{eq:exact-length}
\end{equation}
In particular, $\ell(5)=5$ and $\ell(4093)=15293$.
\end{theorem}

\begin{proof}
Fix $2\le d\le m$.  Because $p$ is prime, the branch begins with
$T_{p,d}\to Q_{p,d,2}$ and then visits the quotient values
\[
2,3,\ldots,\left\lfloor\frac pd\right\rfloor+1.
\]
At the last value $dq>p$, so the edge goes to $T_{p,d+1}$.  From
$T_{p,d}$ through this last quotient state to $T_{p,d+1}$ there are
\[
1+\left\lfloor\frac pd\right\rfloor
=\left\lceil\frac pd\right\rceil
\]
edges; the equality uses $d\nmid p$.  Summing these contributions and adding
$I_p\to T_{p,2}$ plus the contracted terminal return
$T_{p,m+1}\to I_p$ gives \eqref{eq:exact-length}.  Direct substitution gives
the two stated values.
\end{proof}

\begin{corollary}[Bounds and asymptotic]
\label{cor:length-asymptotic}
With $H_m=\sum_{d=1}^m d^{-1}$,
\begin{equation}
p(H_m-1)+2\le\ell(p)\le p(H_m-1)+m+1.
\label{eq:length-bounds}
\end{equation}
The lower inequality is strict for $p\ge5$.  Moreover,
\begin{equation}
\ell(p)=\frac12p\log p+(\gamma-1)p+O(\sqrt p),
\qquad
\frac{\ell(p)}{p\log p}\longrightarrow\frac12.
\label{eq:length-asymptotic}
\end{equation}
\end{corollary}

\begin{proof}
For each tested divisor,
\[
\frac pd\le\left\lceil\frac pd\right\rceil
<\frac pd+1.
\]
There are $m-1$ terms, yielding \eqref{eq:length-bounds}; the upper strict
bound can be weakened to the displayed integer-friendly inequality.  For a
prime and $2\le d\le m$, $p/d$ is not integral, so the lower inequality is
strict whenever this range is nonempty, namely $p\ge5$.

Now $m=\sqrt p+O(1)$ and
$H_m=\log m+\gamma+O(m^{-1})$.  Hence
\[
p(H_m-1)
=\frac12p\log p+(\gamma-1)p+O(\sqrt p),
\]
while the gap in \eqref{eq:length-bounds} is $O(\sqrt p)$.  This proves
\eqref{eq:length-asymptotic}.
\end{proof}

The key quotient is therefore
\begin{equation}
\frac{\log p}{\ell(p)}\sim\frac2p\longrightarrow0.
\label{eq:mean-clock}
\end{equation}
It is the average edge clock forced by the arithmetic total.  No roof
allocation can make its minimum exceed its average, so
\eqref{eq:mean-clock} is already a warning that the expanded computation and
a compact weighted adjacency may be incompatible.
